\section{Memory closures}
\label{sec:memory}

The Mori--Zwanzig formalism \cite{Mori65,Zwanzig61,Zwanzig} projects the
full dynamics onto the retained variables without approximation. It does
not produce a closure. It shows what every closure approximates.

\subsection{The linear case is a Schur complement}

Split a linear system into retained variables $x$ and discarded variables
$y$,
\begin{equation}
  \frac{\dd}{\dd t}\begin{pmatrix}x\\y\end{pmatrix}
  =\begin{pmatrix}A&B\\C&D\end{pmatrix}\begin{pmatrix}x\\y\end{pmatrix}.
\end{equation}
Solving the second row for $y$ and substituting gives
\begin{equation}
  \dot x(t)=A\,x(t)+\int_0^tB\,e^{Ds}C\,x(t-s)\,\dd s+B\,e^{Dt}y(0),
  \label{eq:mz-linear}
\end{equation}
which is \eqref{eq:fastslow-exact} with matrices. Under the Laplace
transform $\hat x(s)=\int_0^\infty e^{-st}x(t)\,\dd t$,
\begin{equation}
  \bigl[s-A-B(s-D)^{-1}C\bigr]\hat x(s)=x(0)+B(s-D)^{-1}y(0) .
  \label{eq:schur}
\end{equation}
The operator on the left is the Schur complement of the resolvent. The
exact closure is the frequency-dependent matrix
$\hat K(s)=B(s-D)^{-1}C$.

Every closure of a linear problem is an approximation of $\hat K(s)$:
\begin{itemize}
  \item \textbf{Quasi-static (Markov):} $\hat K(s)\approx\hat K(0)=-BD^{-1}C$.
    This is the Chapman--Enskog closure. It is accurate when the retained
    dynamics is slow compared with the spectrum of $D$.
  \item \textbf{First correction:}
    $\hat K(s)\approx-BD^{-1}C-s\,BD^{-2}C$, which gives
    $(I+BD^{-2}C)\dot x=(A-BD^{-1}C)x$, the rational closure of
    Example~\ref{ex:fastslow}.
  \item \textbf{Rational in $s$:} approximate $\hat K$ by a ratio of
    polynomials. This is an extended-state model
    (Section~\ref{sec:embedding}).
\end{itemize}

\subsection{The general identity}

For a nonlinear system $\dot X=F(X)$ let $\cL=F\cdot\nabla$ be the
Liouville operator acting on observables, $\cP$ a projection onto functions
of the retained variables, and $\cQ=1-\cP$. The Dyson identity
\[
  e^{t\cL}=e^{t\cQ\cL}+\int_0^te^{(t-s)\cL}\,\cP\cL\,e^{s\cQ\cL}\,\dd s
\]
applied to $\cQ\cL a$ for a retained observable $a$ gives
\begin{equation}
  \frac{\dd}{\dd t}a(t)=e^{t\cL}\cP\cL a
    +\int_0^te^{(t-s)\cL}\,\cP\cL\,e^{s\cQ\cL}\cQ\cL a\,\dd s
    +e^{t\cQ\cL}\cQ\cL a .
  \label{eq:mz}
\end{equation}
The three terms are a Markovian term, a memory term, and the
orthogonal dynamics $F(t)=e^{t\cQ\cL}\cQ\cL a$, which depends on the
discarded initial data and satisfies $\cP F(t)=0$. With Mori's linear
projection onto $a$ in an equilibrium inner product, \eqref{eq:mz} becomes
the generalized Langevin equation
\begin{equation}
  \frac{\dd a}{\dd t}=i\Omega\,a(t)-\int_0^tK(s)\,a(t-s)\,\dd s+F(t),
  \qquad
  K(t)=\avg{F(t),F(0)}\avg{a,a}^{-1} .
  \label{eq:gle}
\end{equation}
The second relation is the fluctuation--dissipation theorem: memory and
noise are two views of the same discarded dynamics, so a closure that keeps
dissipation and drops the noise is correct only for averages.

The identity is exact and not useful as it stands, because the orthogonal
dynamics $e^{t\cQ\cL}$ is as hard as the original problem. What it provides
is a classification.

\subsection{Approximations of the memory term}

\begin{itemize}
  \item \textbf{Short memory.} $K(t)\approx2\gamma\,\delta(t)$. Markovian
    closure. Valid under time-scale separation.
  \item \textbf{Long memory} ($t$-model \cite{ChorinHaldKupferman02}). Take
    the integrand constant in $s$, so the memory term is $t$ times its
    initial value. Suited to problems with no separation.
  \item \textbf{Finite memory.} Truncate the integral at a fixed lag, or
    expand the kernel in a finite basis. The length of the lag is a
    modelling choice that can be tested (hypothesis 3 of the mathematical
    closure program).
  \item \textbf{Data-driven kernels.} Fit $K$ from time series of the
    retained variables \cite{Lin21}.
\end{itemize}

\subsection{Extended states are memory with exponential kernels}
\label{sec:embedding}

If the kernel is a sum of exponentials, $K(t)=\sum_jc_je^{-\lambda_jt}$,
each term can be generated by one auxiliary variable
$\dot z_j=-\lambda_jz_j+c_j\,a$, and the non-Markovian equation for $a$
becomes a Markovian system for $(a,z_1,\dots)$. This is called Markovian
embedding. It runs both ways.

\begin{example}[Cattaneo's law]
\label{ex:cattaneo}
Retaining the heat flux with a relaxation time $\tau$,
\begin{equation}
  \tau\,\partial_tq+q=-\kappa\,\partial_xT
  \qquad\Longleftrightarrow\qquad
  q(t)=-\frac{\kappa}{\tau}\int_0^\infty e^{-s/\tau}\,\partial_xT(t-s)\,\dd s .
\end{equation}
With time dependence $e^{-i\omega t}$ the conductivity is
$\kappa/(1-i\omega\tau)$. Fourier's law is the limit $\omega\tau\to0$. The
temperature equation changes from parabolic to hyperbolic, with signal
speed $\sqrt{\kappa/\tau}$ in suitable units.
\end{example}

Grad's thirteen-moment system is this construction for the Boltzmann
equation: the stress and heat flux are auxiliary variables whose
elimination produces exponential memory kernels in the conservation laws.
Conversely, whether to keep a flux as a state variable is the question of
whether the memory kernel varies over the frequencies the reduced model
excites. The heat-flux manuscript makes this quantitative
(Section~\ref{sec:repo}): in one of its two embedded tests the exact
quasi-static closure has transient errors of up to $34\%$, and in the other
it stays within $2\%$, and the difference is the time scale of the resolved
dynamics.

\subsection{Continued fractions and the moment hierarchy}
\label{sec:cf}

Applying the projection repeatedly, first onto $a$, then onto the part of
$\dot a$ orthogonal to $a$, and so on, generates an orthogonal basis in
which the dynamics is tridiagonal, and expresses the Laplace transform of
the correlation function as a continued fraction (Mori). The construction
is the Lanczos algorithm.

For free streaming, \eqref{eq:linkin} with $\Lambda=0$, this basis is the
Hermite basis and the tridiagonal system is \eqref{eq:hermitehier}. For a
single Fourier mode define $\tau_n(s)$ by
$\hat a_{n}=-ik\sqrt n\,\hat a_{n-1}/\tau_n(s)$ for $n\ge2$. Then
\begin{equation}
  \tau_n(s)=s+\frac{(n+1)\,k^2}{\tau_{n+1}(s)},
  \label{eq:cf}
\end{equation}
and unrolling it gives the classical continued fraction for the Laplace
transform of the free-streaming correlation $e^{-k^2t^2/2}$,
\[
  \int_0^\infty e^{-st}e^{-k^2t^2/2}\,\dd t
  =\cfrac{1}{s+\cfrac{k^2}{s+\cfrac{2k^2}{s+\cfrac{3k^2}{s+\cdots}}}} .
\]
A closure of the moment hierarchy at level $N$ is a choice of
\emph{terminator} for this fraction: a replacement for the exact tail
$\tau_{N+1}(s)$. Truncation takes $\tau_{N+1}=\infty$. Relaxation closures
take $\tau_{N+1}=s+\nu$. The Hammett--Perkins closures take
$\tau_{N+1}$ equal to its exact value at zero frequency, which is
proportional to $|k|$ (Section~\ref{sec:collisionless}). The moment
hierarchy of kinetic theory and the Mori--Zwanzig hierarchy are the same
object.
